\section{0x41\ 数据结构\ 序列结构}
\label{sec:0x41}

\subsection{ST 表}

\InputCode{cpp}{contents/0x40 数据结构/code/0x41 ST 表.cpp}

\subsection{并查集}

\InputCode{cpp}{contents/0x40 数据结构/code/0x41 可撤销并查集.cpp}

\subsection{线段树}

\subsubsection{扫描线}



\subsubsection{势能线段树}

\Highlight{区间开平方}：一个整数 $x$ 开平方 $\mathcal{O}(\log \log V)$ 次之后就会变成 $1$。若区间 $\mathrm{max} = 1$ 则 \texttt{return}。

\Highlight{区间取模}：一个整数 $x$ ($x \ge p$) 对 $p$ 取模，$x$ 至少变小一半。若区间 $\mathrm{max} < p$ 则 \texttt{return}。

\Highlight{区间除法}：整除 $p$ 之后，由于 $\lfloor \frac{\max}{p} \rfloor - \lfloor \frac{\min}{p} \rfloor \leq \lceil \frac{\max - \min}{p} \rceil$，区间最大值与最小值的差，多数情况下会减半。

\begin{itemize}
\item 若区间 $\mathrm{max} - \mathrm{min} = 0$，则区间加之后 \texttt{return}。
\item 若区间 $\mathrm{max} - \mathrm{min} = 1$ 且 $\lfloor \frac{\mathrm{\max}}{p} \rfloor \neq \lfloor \frac{\mathrm{\min}}{p} \rfloor$，则区间加之后 \texttt{return}。
\end{itemize}

\Highlight{区间按位与}：一个二进制数 $x$ 与 $v$ 进行按位与，至多进行 $\log V$ 次有效操作。若区间（在二进制下）$\mathrm{bitor} \subseteq v$ 则 \texttt{return}。

\Highlight{区间取最值（吉老师线段树）}：以区间取最小值为例：

\begin{itemize}
\item 在线段树上的每一个节点 $p$ 需要维护：
  \begin{itemize}
  \item \texttt{mx}：区间最大值。
  \item \texttt{cnt}：区间最大值个数。
  \item \texttt{md}：区间严格次大值。
  \end{itemize}
\item 在区间修改时，若遇到一个 “被询问区间完全覆盖的节点” 时：
  \begin{itemize}
  \item 若 $x \geq \mathrm{mx}$，则 \texttt{return}。
  \item 若 $x \leq \mathrm{md}$，则暴力递归左右儿子修改。
  \item 若 $\mathrm{md} < x < \mathrm{mx}$，则有 $\mathrm{cnt}$ 个 $\mathrm{mx}$ 被修改成 $x$，打标记后 \texttt{return}。
  \end{itemize}
\item 时间复杂度 $\mathcal{O}(n \log^2 n)$。
\end{itemize}

\InputCode{cpp}{contents/0x40 数据结构/code/0x41 吉老师线段树.cpp}

\subsection{笛卡尔树}



\subsubsection{区间笛卡尔树}

